\subsection{Two Systems, or One? Dual-Process Theories of Moral Judgment}
\label{sec:5.3.1}
The framework organizing most discussion of moral judgment is dual-process theory, and though it reached wide audiences through Daniel Kahneman's account of a fast, intuitive "System 1" and a slow, deliberate "System 2," the idea was built over decades by many researchers — Peter Wason and Jonathan Evans, Seymour Epstein, Steven Sloman, and Keith Stanovich among them \autocite{evans2013dual}. The two systems, in the standard telling, have opposite characters. The first is fast, automatic, parallel, effortless, and emotionally colored — the source of snap judgments and gut feelings. The second is slow, serial, effortful, and rule-governed — what a person uses to check a sum, follow an argument, or resist a tempting conclusion. The first runs constantly and almost for free. The second is expensive to run and easily left idle.

The framework's reach into moral judgment specifically runs through the psychologist Joshua Greene, who used it to explain why ethical intuitions contradict themselves. Consider the trolley problem in its two classic forms. A runaway trolley will kill five people unless diverted onto a side track, where it will kill one; most people say divert it. Now stop the same trolley by pushing a stranger off a footbridge into its path — same arithmetic, five saved for one — and most people refuse. Greene's brain-imaging work found a difference between the two cases in exactly the region the theory would predict: the up-close, hands-on violence of the footbridge case engaged emotional-processing regions that stayed quiet during the impersonal switch case, and that difference tracked which choice people made \autocite{greene2001fmri}. A separate line of evidence sharpens the same point. Patients with damage to the ventromedial prefrontal cortex, a region central to normal emotional response, endorse the footbridge push at a higher rate than people with intact emotional processing do — specifically on the personal, high-stakes dilemmas alone \autocite{koenigs2007damage}. Remove the emotional alarm, and the arithmetic wins more often.

All of which makes dual-process theory sound more settled than it is. Its critics — including Evans and Stanovich, the framework's own leading defenders — argue that the cartoon version, two little agents in the head taking turns at the controls, explains everything after the fact and predicts nothing beforehand: almost any result can be narrated as System 1 overruled by System 2, or System 2 captured by System 1. Evans and Stanovich have spent years trying to rescue the framework from its own popularity, insisting there are not two systems but two types of processing — one fast and autonomous, one slow and demanding of working memory — and that the tidy personification is a teaching device and makes no claim about the brain's architecture \autocite{evans2013dual}. Whether "dual-process" names two mechanisms or one continuum under different load is not resolved.

That unresolved status matters for how an AI system's moral judgment gets built, and it is a design question before it is a descriptive one. A designer who takes the cartoon version literally is tempted toward a literal architecture: a fast heuristic layer bolted to a slow deliberative one, with the deliberative layer assumed to be the trustworthy check on the heuristic one. The trolley evidence complicates that assumption before anything else does — the "cooler" answer produced by weaker emotional engagement is not obviously the more correct one. It is simply the one with less flinch in it. A system built to override its fast pattern-matching with slower rule-following whenever the two disagree is not thereby built to be more moral; it is built to resemble the profile of the rare human patients for whom losing the flinch is a symptom, not an achievement.
